%%%%%%%%%%%%%%%%%%%%%%%%%%%%%%%%%%%%%%%%%%%%%%%%%%%%%%%%%%%%%%%%%%%%%%%%%%%%%%%%%%%%%%%%%%%%%%%%%%%%%%%%%%%%%%%%%%%%%%%%%%%%%%%%%%%%%%%%%%%%%%%%%%%%%%%%%%%%%%%%%%%%%%%%%%%%%%%%%%%%%%%%%%%%%%%%%%%%%%%%%%%%%%%%%%%%%%%%%%%%%%%%%%%%%%%%%%%%%%%%%%%%%%%%%%%%
\documentclass[12pt, twosides]{article}
\usepackage{graphicx}
\usepackage{cite}
\usepackage{picinpar}
\usepackage{amsmath}
\usepackage{url}
\usepackage{flushend}
\usepackage[latin1]{inputenc}
\usepackage{colortbl}
\usepackage{soul}
\usepackage{multirow}
\usepackage{amsfonts}
\usepackage{amssymb}
\usepackage{pifont}
\usepackage{booktabs}
\usepackage{color}
\usepackage{alltt}
\usepackage[hidelinks]{hyperref}
\usepackage{enumerate}
\usepackage{siunitx}
\usepackage{breakurl}
\usepackage{caption}
\usepackage{epstopdf}
\usepackage{pbox}
\usepackage{stfloats}
\usepackage{algorithm}
\usepackage{algorithmic}
\usepackage{changebar}
\newcommand{\vect}[1]{\mathbf{#1}}
\usepackage{amsthm}
\usepackage{diagbox,tabu,enumerate,fancybox}
\newtheorem{theorem}{Theorem}
\newtheorem{lemma}{Lemma}
\newtheorem{remark}{Remark}
\long\def\comment#1{ }

\usepackage[top=1in, bottom=1in, left=1.2in, right=1.2in]{geometry}
\setlength{\parskip}{1em}

\theoremstyle{definition}
\newtheorem{defn}{Definition}[section]
\newtheorem{conj}{Conjecture}[section]
\newtheorem{exmp}{Example}[section]
\newenvironment{fminipage}%
{\begin{Sbox}\begin{minipage}}%
{\end{minipage}\end{Sbox}\doublebox{\TheSbox}}

\begin{document}
\bigskip

\noindent \textbf{Title:Event-Triggered Nonlinear Model Predictive Control for Energy-Aware UAV Obstacle Avoidance}

\noindent \textbf{Authors: Heng Zhang, ~Liang Yuan, ~Xianghui Cao}

\noindent \textbf{Reference:  IoT-69021-2026}

\bigskip

\begin{center}
\textbf{\LARGE Authors' Responses to the Reviewers}
\end{center}

\bigskip First we would like to express our sincere thanks to the Associate Editor and the anonymous reviewers for their valuable comments and suggestions on our manuscript. The following is a detailed description of the reviewer comments and the currently completed response.

%%%%%%%%%%%%%%%%%%%%%%%%%%%%%%%%%%%%%%%%%%%%%%%%%%%%%%%%%%%%%%
\begin{center}
\textbf{\textit{\large Response to the Associate Editor}}
\end{center}

\textcolor[rgb]{0.00,0.00,1.00}{[\textbf {Associate Editor}]} \textit{\textbf{Based on 3 qualified reviews, this AE makes Major Revision decision. The main concerns include weak connection to the Internet of Things, unclear real-time feasibility claims, unclear mechanism that drives the UAV toward the target during NMPC replanning,  and so on.}}

\textcolor[rgb]{0.00,0.00,1.00}{[\textbf{Authors}]}

%%%%%%%%%%%%%%%%%%%%%%%%%%%%%%%%%%%%%%%%%%%%%%%%%%%%%%%%%%%%%%
\begin{center}
\textbf{\textit{\large Response to Reviewer \# 1}}
\end{center}

\textcolor[rgb]{0.00,0.00,1.00}{[\textbf {Reviewer}]} \textit{\textbf{Comments to the Author}}

\textit{\textbf{This manuscript proposes an event-triggered nonlinear model predictive control framework for energy-aware UAV obstacle avoidance in complex three-dimensional dynamic environments. The framework combines obstacle-motion prediction, composite risk assessment, dual-mode control, and an asymmetric climb penalty. The manuscript is generally well organized, and the authors have conducted comparative simulations, parameter-sensitivity analysis, computational profiling, and Monte Carlo experiments.}}

\textit{\textbf{The topic is relevant to autonomous UAV systems and may be of interest for computationally constrained applications. However, several issues regarding the IoT-oriented positioning of the work, real-time feasibility, optimization formulation, validation of the claimed contributions, and experimental methodology need to be addressed before the manuscript can be considered for publication.}}

\textcolor[rgb]{0.00,0.00,1.00}{[\textbf{Authors}]}

\textit{\textbf{Major Comments}}

\textcolor[rgb]{0.00,0.00,1.00}{[\textbf {Reviewer}]} \textit{\textbf{1. The connection to the Internet of Things should be strengthened.}}

\textit{\textbf{The current manuscript is primarily framed as a UAV control and dynamic obstacle-avoidance study. Although autonomous UAVs are relevant to many IoT applications, the IoT context remains somewhat peripheral to the technical development in the present version.}}

\textit{\textbf{The manuscript would benefit from a clearer description of the UAV's role within an IoT or autonomous IoT architecture. In particular, it should be clarified how the event-triggered mechanism translates into resource savings from an IoT-system perspective, for example, reduced onboard computational load or edge-computing demand. If communication overhead, sensing-data exchange, or other network-related resource consumption is also expected to benefit from the proposed framework, such effects should be explicitly modeled or discussed rather than only stated conceptually.}}

\textit{\textbf{I suggest revising the introduction, system description, and conclusion so that the IoT motivation is consistently connected to the proposed control architecture and its resource-saving mechanism.}}

\textcolor[rgb]{0.00,0.00,1.00}{[\textbf{Authors}]}

Thank you for this valuable comment. We agree that, in the original manuscript, the connection between the proposed framework and the IoT system context was not sufficiently emphasized. Although the study considered UAV-assisted obstacle avoidance and energy-aware control, the role of the UAV and the proposed control framework within an IoT-enabled cyber-physical system required clearer system-level explanation.

In response, we have revised the Abstract, Introduction, Section II, and Conclusion to clarify the IoT relevance and scope of this work. In the revised manuscript, the UAV is explicitly described as a mobile sensing-and-actuation node in an IoT-enabled cyber-physical system, where it operates through a local perception-decision-actuation loop. The revised system description further explains that obstacle observations are processed locally for prediction, risk assessment, mode selection, and control action generation. In particular, $y_i(k)$ is clarified as the available obstacle observation provided to the local perception-control loop.

We have also clarified the role of the event-triggered mechanism from the perspective of resource-aware local computation. Specifically, the mechanism avoids unnecessary NMPC solver activations during low-risk conditions and reduces the associated local controller-computation workload, while nominal tracking control remains available under low-risk conditions. Therefore, the IoT relevance of the proposed approach is positioned as a resource-aware local perception-control component for an IoT-enabled cyber-physical system.

To avoid overstatement, we explicitly clarify that communication scheduling, radio resource allocation, and communication-energy optimization are outside the modeling scope of this work. Accordingly, the revised manuscript does not claim communication optimization or network-resource benefits; rather, it focuses on the local perception-control and computational decision process of the UAV node.

The corresponding modifications are shown below:

"The framework is positioned in an IoT-enabled cyber--physical system, where the UAV operates as a mobile sensing-and-actuation node that converts available obstacle observations into safety-critical motion decisions."

"In an IoT-enabled cyber--physical system, the UAV can operate as a mobile sensing-and-actuation node that closes a local perception--decision--actuation loop: it uses environmental observations for local situation assessment, selects a control action, and executes the resulting maneuver."

"Within this context, the present work focuses on the local perception--control computation required for dynamic obstacle avoidance. It does not model communication scheduling, radio resource allocation, or communication-energy optimization; these aspects are outside the scope of the proposed control formulation."

"From an IoT-system perspective, the event-triggered mechanism performs resource-aware local decision scheduling by avoiding unnecessary NMPC solver activations and reducing the associated local controller-computation workload."

"The UAV is modeled as a mobile sensing-and-actuation node in an IoT-enabled cyber--physical system. At each sampling instant, the local perception--decision--actuation loop processes available obstacle observations, predicts their future motion, assesses collision risk, selects an operating mode, and applies the corresponding control action. This work models this local loop and its motion consequences."

"Here, $\vect{y}_i(k)$ represents the available obstacle observation provided to the local perception--control loop at time step $k$."

"From an IoT-system perspective, the proposed method provides a resource-aware local perception--control component for a mobile sensing-and-actuation node, with the demonstrated benefit of avoiding unnecessary NMPC activations and associated local computational workload."

The corresponding modifications have been highlighted in blue in the revised manuscript.

\medskip

\textcolor[rgb]{0.00,0.00,1.00}{[\textbf {Reviewer}]} \textit{\textbf{2. The real-time feasibility claims need to be reconsidered.}}

\textit{\textbf{The manuscript uses a sampling interval of 0.1 s, whereas the computational profiling results report an average CPU time per step of approximately 2074 ms and an average CPU time of approximately 2658 ms for triggered steps. These values are far larger than the stated 100-ms control interval.}}

\textit{\textbf{Therefore, the current MATLAB/CasADi implementation does not appear to satisfy the stated closed-loop update frequency. The manuscript should clarify how simulation time is advanced relative to solver wall-clock time and whether computation delay is incorporated into the closed-loop dynamics. Given the reported timing results, it is also important to distinguish computational savings relative to continuous NMPC from actual real-time executability.}}

\textit{\textbf{The reduction in unnecessary NMPC calls is still a meaningful computational advantage. However, reducing average computational cost relative to continuous NMPC is not equivalent to meeting a real-time deadline. Accordingly, the real-time feasibility claims in the abstract, introduction, Section IV-D, and conclusion should be moderated.}}

\textit{\textbf{A compiled implementation, hardware-in-the-loop experiment, or onboard/embedded validation would substantially strengthen this aspect of the work. At minimum, I suggest reporting deadline-oriented timing statistics, including the median, 95th or 99th percentile computation time, and the percentage of control steps exceeding the 0.1-s sampling deadline.}}

\textcolor[rgb]{0.00,0.00,1.00}{[\textbf{Authors}]} Thank you very much for the reviewer's insightful comment regarding the inconsistency between the real-time feasibility claims and the reported computational times. We appreciate the reviewer's request for a clearer and more rigorous disclosure of the computational scope and deadline-related limitations.

We agree that the term ``real-time feasibility'' was too broad in the original manuscript. The reported experiments are based on desktop MATLAB/CasADi profiling, rather than an embedded deployment, and the measured CPU times do not provide a hard real-time guarantee at the simulated sampling interval of $\Delta t=0.1$ s. Although event triggering reduces unnecessary NMPC activations and achieves relative computational savings compared with continuous NMPC, this reduction is not equivalent to demonstrating that all control updates satisfy a 100-ms deadline.

In the revised manuscript, we have moderated the computational claims in four locations. First, the Abstract now describes computational efficiency and resource-aware computation reduction, and identifies the reported 21.1\%--24.2\% CPU-time reduction as a desktop MATLAB/CasADi profiling result. Second, the Introduction now explicitly states that the present study provides desktop computational characterization rather than strict onboard real-time validation. Third, Section IV-D has been retitled ``Computational Performance Profiling and Deadline-Oriented Analysis.'' It now specifies the timing scope, reports deadline-oriented statistics, and distinguishes desktop controller-computation time from complete physical closed-loop latency. Finally, the Conclusion now characterizes the results as desktop computational behavior and relative CPU savings, while identifying embedded platforms, hardware-in-the-loop experiments, and compiled implementations as future validation work.

The corresponding modifications are shown below:

"computational efficiency and resource-aware computation reduction"

"Compared with the continuous NMPC baseline, the proposed method reduces average CPU time by 21.1\%--24.2\% in the reported desktop MATLAB/CasADi profiling experiments while preserving reliable clearance and stable solver behavior."

"The event-triggered mechanism reduces unnecessary optimization activations and computational workload, while the current study focuses on desktop MATLAB/CasADi computational characterization rather than strict onboard real-time validation. Accordingly, the reported CPU-time reductions represent relative computational savings, not a hard real-time guarantee."

"The reported computation time measures the controller computation process in MATLAB/CasADi simulation, including the triggered NMPC solving procedure, rather than the complete physical closed-loop latency including communication and hardware execution."

"Table V also reports the IPOPT statistics and solver configuration for the reported desktop MATLAB/CasADi implementation. The average IPOPT iteration number remains below the maximum iteration limit, and no solver failure or fallback occurs during the profiling runs. The profiling configuration uses IPOPT with the MUMPS sparse linear solver, fixed convergence tolerances, and shifted-reference warm start. The memory values refer to the MATLAB process during the reported runs; they indicate no continuous process-memory growth in this desktop experiment, but they are not solver-specific or embedded-memory measurements."

"Although the event-triggered strategy significantly reduces unnecessary NMPC calls, the current desktop MATLAB/CasADi implementation does not satisfy the strict 100-ms deadline for triggered optimization steps. With $\Delta t=0.1$ s, the deadline-violation ratio is 78.168\% over all recorded controller steps and 100.0\% over the triggered NMPC steps. Therefore, the reported results characterize desktop computational behavior and relative computational savings rather than hard real-time or onboard feasibility. Future validation on embedded platforms, hardware-in-the-loop experiments, and compiled implementations will be considered in future work."

"Computational profiling also reported IPOPT iterations, solver settings, MATLAB process-memory observations, and timing distributions. These results characterize desktop computational behavior and relative CPU savings; they do not establish hard real-time operation at the 0.1-s deadline or embedded feasibility."

"Future work will consider validation on embedded platforms, hardware-in-the-loop experiments, and compiled implementations, together with prediction enhancement and theoretical extension. We will extend the constant-velocity prediction model to more nonlinear and uncertain obstacle motions, study adaptive weighting mechanisms and terminal-set-based stability analysis, and investigate cooperative multi-UAV obstacle avoidance in real cluttered environments."

The complete computational-profiling table has also been revised as follows:

\begin{table*}[t]
	\captionsetup{labelformat=empty}
	\caption{Table IV: Deadline-Oriented Computational Statistics and Solver Disclosure of the Proposed ET-NMPC}
	\label{tab:response_computational_profiling}
	\centering
	\footnotesize
	\setlength{\tabcolsep}{2.5pt}
	\renewcommand{\arraystretch}{1.12}
	\begin{tabular}{p{0.46\textwidth}p{0.46\textwidth}}
		\toprule
		\textbf{Item} & \textbf{Value} \\
		\midrule
		\multicolumn{2}{l}{\textbf{Profiling setup}} \\
		Number of profiling runs & 30 \\
		Total simulation steps & 4072 \\
		Number of dynamic obstacles & 5 \\
		Triggered steps & 3183 \\
		Trigger rate & $77.8 \pm 3.5\%$ \\
		\midrule
		\multicolumn{2}{l}{\textbf{IPOPT statistics}} \\
		Average IPOPT iterations & $22.0 \pm 12.9$ \\
		Maximum IPOPT iterations & 87 \\
		Solver failure count & 0 \\
		Fallback count & 0 \\
		\midrule
		\multicolumn{2}{l}{\textbf{Memory statistics}} \\
		Mean memory before simulation & 6361.1 MB \\
		Mean memory after simulation & 6356.9 MB \\
		Mean memory variation & -4.2 MB \\
		Peak observed memory & 6367.6 MB \\
		\midrule
		\multicolumn{2}{l}{\textbf{Solver settings}} \\
		IPOPT maximum iterations & 100 \\
		IPOPT tolerance & $10^{-4}$ \\
		IPOPT acceptable tolerance & $10^{-3}$ \\
		Sparse linear solver & MUMPS \\
		Warm-start strategy & Shifted reference initialization \\
		\bottomrule
	\end{tabular}

	\medskip
	\textcolor{blue}{\textbf{Deadline-oriented timing statistics (pooled step-level values)}}\\
	\resizebox{\textwidth}{!}{%
		\begin{tabular}{lccccc}
			\toprule
			\textcolor{blue}{\textbf{Metric}} & \textcolor{blue}{\textbf{Mean (ms)}} & \textcolor{blue}{\textbf{Median (ms)}} & \textcolor{blue}{\textbf{95th percentile (ms)}} & \textcolor{blue}{\textbf{99th percentile (ms)}} & \textcolor{blue}{\textbf{Maximum (ms)}} \\
			\midrule
			\textcolor{blue}{All controller steps} & \textcolor{blue}{2077.9} & \textcolor{blue}{2551.9} & \textcolor{blue}{2964.7} & \textcolor{blue}{3150.4} & \textcolor{blue}{3982.8} \\
			\textcolor{blue}{Triggered NMPC steps} & \textcolor{blue}{2658.3} & \textcolor{blue}{2596.2} & \textcolor{blue}{2987.2} & \textcolor{blue}{3187.8} & \textcolor{blue}{3982.8} \\
			\midrule
			\textcolor{blue}{Deadline} & \multicolumn{5}{c}{\textcolor{blue}{100 ms ($\Delta t=0.1$ s)}} \\
			\textcolor{blue}{Violation ratio (all steps)} & \multicolumn{5}{c}{\textcolor{blue}{3183/4072 (78.168\%)}} \\
			\textcolor{blue}{Violation ratio (triggered steps)} & \multicolumn{5}{c}{\textcolor{blue}{3183/3183 (100.0\%)}} \\
			\bottomrule
		\end{tabular}%
	}
\end{table*}

\clearpage

These revisions more accurately distinguish the demonstrated computational efficiency and relative CPU savings from hard real-time capability. The corresponding modifications have been highlighted in blue in the revised manuscript.

\medskip

\textcolor[rgb]{0.00,0.00,1.00}{[\textbf {Reviewer}]} \textit{\textbf{3. The mechanism that drives the UAV toward the target during NMPC replanning is unclear.}}

\textit{\textbf{The objective function in (19)--(24) includes path-length, smoothness, energy, and collision-risk terms. However, I do not see an explicit reference-tracking term, path-progress term, or terminal goal cost in the high-risk NMPC formulation.}}

\textit{\textbf{As currently written, it is unclear what mechanism promotes continued progress toward the target while the predictive optimizer is active, or how the UAV is encouraged to return to the nominal trajectory after an avoidance maneuver. In particular, because path length, control effort, and energy consumption are all penalized, a solution with limited motion could potentially be favored once the safety constraints are satisfied if progress toward the target is not explicitly rewarded.}}

\textit{\textbf{If a reference-tracking cost, path-progress term, terminal goal cost, or equivalent mechanism is used in the actual implementation, it should be explicitly included in the mathematical formulation. Otherwise, the objective formulation may need to be revised, and the corresponding simulation results should be updated accordingly.}}

\textcolor[rgb]{0.00,0.00,1.00}{[\textbf{Authors}]} Thank you very much for the reviewer's important comment regarding the target-oriented mechanism of the high-risk NMPC formulation.

We agree that the original manuscript objective function lacked explicit reference-tracking, path-progress, and terminal-attraction mechanisms. Although the original formulation included path-length, smoothness, energy, and collision-risk terms, it did not explicitly state how the optimizer should remain aligned with the local reference or how it should be attracted toward the terminal target after an avoidance maneuver. Consequently, the target-oriented behavior during replanning was not sufficiently clear in the original manuscript.

In the revised manuscript, Section III-C has been revised to make this mechanism explicit. First, a reference-tracking cost has been introduced to penalize the deviation between each predicted UAV position and the corresponding local reference position, thereby preserving reference-following behavior during local replanning. Second, a terminal goal cost has been added to penalize terminal position error and terminal velocity. Third, a stage-wise terminal-attraction term has been incorporated into the terminal cost. Its influence increases over the prediction horizon, so that the predicted trajectory is progressively attracted toward the conditional terminal reference. Together, these terms explicitly connect local obstacle avoidance with continued progression toward the target.

The revised objective function combines the reference-tracking, control-increment, energy-aware, differentiable soft collision-avoidance, and terminal costs as follows:
\[
J = \rho_1 J_{\mathrm{track}} + \rho_2 J_{\mathrm{smooth}} + \rho_3 J_{\mathrm{energy}} + \rho_4 J_{\mathrm{risk}} + \rho_5 J_{\mathrm{terminal}}.
\]

The corresponding objective weights are $\rho_1=1.0$, $\rho_2=0.5$, $\rho_3=0.8$, $\rho_4=5.0$, and $\rho_5=1.0$. The terminal-cost weights are $\lambda_p=120$, $\lambda_v=20$, and $\lambda_{\mathrm{prog}}=8$. These parameters are reported in the revised Table II and are consistent with the MATLAB/CasADi NMPC solver.

The corresponding modifications are shown below:

"To balance local-reference tracking, control regularization, energy awareness, and differentiable soft collision avoidance, the overall objective function $J$ is defined as
\[
J = \rho_1 J_{\mathrm{track}} + \rho_2 J_{\mathrm{smooth}} + \rho_3 J_{\mathrm{energy}} + \rho_4 J_{\mathrm{risk}} + \rho_5 J_{\mathrm{terminal}},
\]
where $J_{\text{track}}$, $J_{\text{smooth}}$, $J_{\text{energy}}$, $J_{\text{risk}}$, and $J_{\text{terminal}}$ denote the reference-tracking, control-increment, energy-aware, soft collision-avoidance, and terminal costs, respectively; $\rho_1,\ldots,\rho_5$ are their nonnegative weights. The risk-fusion weights $w_1,w_2,w_3$ retain their distinct meaning in Section~III-B.

The reference-tracking cost is defined as
\[
J_{\mathrm{track}} = \sum_{\tau=0}^{H-1} \left\|\vect{p}_{k+\tau+1}-\vect{r}_{k+\tau+1}\right\|_2^2,
\]
where $\vect{r}_{k+\tau+1}$ is the local reference position at the corresponding prediction step. This term ensures that the UAV tracks the local reference trajectory during optimization. The control-increment penalty is defined as
\[
J_{\mathrm{smooth}} = \sum_{\tau=0}^{H-2} \left\|\vect{u}_{k+\tau+1}-\vect{u}_{k+\tau}\right\|_2^2,
\]
which penalizes large variations in the predicted control sequence to improve trajectory executability.

To explicitly account for terminal-goal attainment and progressively stronger attraction over the prediction horizon, the terminal cost is defined as
\[
\begin{aligned}
J_{\mathrm{terminal}} ={}& \lambda_p\left\|\vect{p}_{k+H}-\vect{p}_{\mathrm{terminal\_ref}}\right\|_2^2+\lambda_v\left\|\vect{v}_{k+H}\right\|_2^2 \\
&+\lambda_{\mathrm{prog}}\sum_{\tau=0}^{H-1}\left(\frac{\tau+1}{H}\right)^2\left\|\vect{p}_{k+\tau+1}-\vect{p}_{\mathrm{terminal\_ref}}\right\|_2^2,
\end{aligned}
\]
where the first two terms penalize terminal position error and terminal velocity, respectively. The final term is a stage-wise terminal-attraction penalty whose influence increases over the horizon. The terminal reference is $\vect{p}_{\text{goal}}$ when $\left\|\vect{p}_{k}-\vect{p}_{\text{goal}}\right\|_2\leq25\,\mathrm{m}$; otherwise, it is the last local reference point. Minimizing $J_{\text{terminal}}$ therefore promotes terminal convergence while avoiding residual terminal motion."

The corresponding parameters in Table II have also been updated as follows:

\begin{table}[h!]
\centering
\captionsetup{labelformat=empty}
\caption{Table II: Key Parameters Used in the Simulations}
\label{tab:response_key_parameters}
\scriptsize
\renewcommand{\arraystretch}{1.12}
\setlength{\tabcolsep}{2.5pt}
\resizebox{\columnwidth}{!}{%
\begin{tabular}{l l c c}
\toprule
\textbf{Symbol} & \textbf{Description} & \textbf{Unit} & \textbf{Value} \\
\midrule
\multicolumn{4}{l}{\textit{UAV kinematic parameters}} \\
\midrule
$v_{\max}$ & Maximum UAV speed & m/s & 5.0 \\
$u_{\max}$ & Maximum UAV acceleration & m/s$^2$ & 3.0 \\
$z_{\min}$ & Minimum flight altitude & m & 0.5 \\
$z_{\max}$ & Maximum flight altitude & m & 20.0 \\
$r_u$ & UAV safety radius & m & 0.3 \\
$\Delta t$ & Sampling time & s & 0.1 \\
$H$ & Prediction horizon & steps & 15 \\
\midrule
\multicolumn{4}{l}{\textit{NMPC cost weights}} \\
\midrule
\textcolor{blue}{$\rho_1$} & \textcolor{blue}{Reference-tracking cost weight} & \textcolor{blue}{--} & \textcolor{blue}{1.0} \\
\textcolor{blue}{$\rho_2$} & \textcolor{blue}{Control-increment cost weight} & \textcolor{blue}{--} & \textcolor{blue}{0.5} \\
\textcolor{blue}{$\rho_3$} & \textcolor{blue}{Energy-aware cost weight} & \textcolor{blue}{--} & \textcolor{blue}{0.8} \\
\textcolor{blue}{$\rho_4$} & \textcolor{blue}{Differentiable soft risk-cost weight} & \textcolor{blue}{--} & \textcolor{blue}{5.0} \\
\textcolor{blue}{$\rho_5$} & \textcolor{blue}{Outer terminal-cost weight} & \textcolor{blue}{--} & \textcolor{blue}{1.0} \\
$c_1$ & Distance-related energy weight & -- & 1.0 \\
$c_2$ & Control-effort energy weight & -- & 0.1 \\
$c_3$ & Asymmetric climb penalty weight & -- & 3.0 \\
\textcolor{blue}{$\lambda_p$} & \textcolor{blue}{Terminal position weight} & \textcolor{blue}{--} & \textcolor{blue}{120} \\
\textcolor{blue}{$\lambda_v$} & \textcolor{blue}{Terminal velocity weight} & \textcolor{blue}{--} & \textcolor{blue}{20} \\
\textcolor{blue}{$\lambda_{\mathrm{prog}}$} & \textcolor{blue}{Stage-wise terminal-attraction/path weight} & \textcolor{blue}{--} & \textcolor{blue}{8} \\
\midrule
\multicolumn{4}{l}{\textit{Risk assessment and event-triggering parameters}} \\
\midrule
$w_1$ & TCPA-based encounter risk weight & -- & 0.4 \\
$w_2$ & Predictive occupancy risk weight & -- & 0.4 \\
$w_3$ & Reference deviation risk weight & -- & 0.2 \\
$R_{on}$ & Risk activation threshold & -- & 0.6 \\
$R_{off}$ & Risk deactivation threshold & -- & 0.2 \\
$T_{ref}$ & Refractory interval & steps & 5 \\
\bottomrule
\end{tabular}%
}
\end{table}

The corresponding modifications have been highlighted in blue in the revised manuscript.

\medskip

\textcolor[rgb]{0.00,0.00,1.00}{[\textbf {Reviewer}]} \textit{\textbf{4. The claimed benefit of the asymmetric climb penalty is not sufficiently supported by the reported results.}}

\textit{\textbf{The asymmetric climb penalty is presented as one of the main contributions of the manuscript. However, the results in Table III do not clearly demonstrate that the proposed formulation suppresses unnecessary climbing or reduces energy consumption. In fact, the symmetric-energy variant reports slightly lower climb and estimated energy consumption than the proposed method in that table.}}

\textit{\textbf{The current evidence therefore does not appear sufficient to support a strong claim that the asymmetric penalty directly reduces unnecessary vertical maneuvers or energy consumption. More precisely, the additional term in (22) penalizes positive altitude gain rather than vertical motion in general, and its specific contribution should be demonstrated more clearly.}}

\textit{\textbf{A focused ablation study would be useful. For example, the proposed asymmetric penalty could be compared with the symmetric-energy formulation in scenarios where both horizontal and vertical detours are feasible and where climbing has a meaningful energetic consequence. Relevant metrics could include the vertical trajectory, accumulated positive altitude gain, flight time, path length, and estimated energy consumption over multiple randomized trials.}}

\textit{\textbf{Statistical variation should also be reported to determine whether the observed differences are meaningful. If the main effect of the asymmetric penalty is to alter the safety-energy trade-off rather than consistently reduce climb or total energy consumption, the contribution should be described in those more precise terms.}}

\textcolor[rgb]{0.00,0.00,1.00}{[\textbf{Authors}]}

Thank you for this important comment. We agree that the original manuscript did not sufficiently isolate the contribution of the asymmetric climb penalty. Although this penalty was presented as a main contribution, Table III reports an overall framework comparison rather than an ablation that isolates the $c_3$ term. Therefore, the original results could not fully support a strong conclusion about the independent influence of the asymmetric penalty on cumulative positive climb or the estimated energy proxy.

In response, we have added a 30-trial paired Monte Carlo ablation study, entitled \textit{Paired Ablation of the Asymmetric Climb Penalty}, in Section IV, together with a new Table IV. We have also revised the Abstract, Introduction, energy-aware cost formulation, baseline discussion, parameter-sensitivity discussion, and Conclusion. These revisions avoid overstatement and describe the contribution in terms of energy-aware altitude regulation, optimization-preference shaping, and the energy-related trade-off.

The new ablation study compares Asymmetric ET-NMPC with Symmetric-energy ET-NMPC. For each paired trial, both controllers use identical random seeds, identical static maps, identical UAV initial states, identical goal positions, identical dynamic-obstacle realizations, and identical measurement-noise sequences. The only changed setting is the asymmetric positive-climb penalty: Asymmetric ET-NMPC uses $c_3=3.0$, whereas Symmetric-energy ET-NMPC sets $c_3=0$. All other controller, prediction, risk-assessment, event-triggering, and scenario parameters remain unchanged.

The paired results show that Asymmetric ET-NMPC achieves 30/30 successful trials with 0/30 collisions, whereas Symmetric-energy ET-NMPC achieves 29/30 successful trials with 1/30 collision. With paired differences defined as Asymmetric $-$ Symmetric, the total positive-climb difference is 0.217 m with a 95\% confidence interval of [$-0.069$, $0.503$], and the estimated energy-proxy difference is 0.00368 Wh with a 95\% confidence interval of [$-0.00169$, $0.00904$]. Both confidence intervals include zero. Therefore, the experiment does not provide statistical evidence that the asymmetric penalty reduces cumulative positive climb or estimated energy proxy under the tested parameterization.

Accordingly, the asymmetric climb penalty is retained as an energy-aware altitude regulation mechanism, an optimization-preference shaping mechanism, and a component for regulating the safety--energy trade-off. It assigns additional cost to positive altitude gains without claiming a guaranteed lower energy outcome or suppression of unnecessary climbing. The estimated energy is explicitly described as a model-based mission-energy proxy for comparative analysis, rather than a hardware battery measurement.

The corresponding modifications are shown below:

"C. Paired Ablation of the Asymmetric Climb Penalty"

"To isolate the influence of the asymmetric climb penalty, we conducted 30 paired Monte Carlo scenarios comparing Asymmetric ET-NMPC with Symmetric-energy ET-NMPC. For each paired scenario, both controllers used the same random seed, static map, UAV initial state, goal position, dynamic-obstacle realization, and measurement-noise sequence. The only changed setting was the additional positive-climb penalty: the asymmetric controller used $c_3=3.0$, whereas the symmetric-energy controller disabled this term by setting $c_3=0$."

"The asymmetric controller achieved 30/30 successful trials with no collision, whereas the symmetric-energy controller achieved 29/30 successful trials with one collision. The paired difference in total positive climb was 0.217 m (95\% CI: [$-0.069$, $0.503$]), and the paired difference in estimated energy proxy was 0.00368 Wh (95\% CI: [$-0.00169$, $0.00904$]), where each difference is defined as Asymmetric $-$ Symmetric. Both confidence intervals include zero. Therefore, this experiment does not provide statistical evidence that the asymmetric penalty reduces cumulative positive climb or estimated mission energy under the tested parameterization. Instead, the additional term shapes altitude-related optimization preference and provides energy-aware altitude regulation within the safety--energy tradeoff."

"An asymmetric climb penalty is also introduced as an energy-aware altitude regulation mechanism that shapes the optimization preference for positive altitude gains."

"These designs reduce unnecessary NMPC calls, schedule replanning according to real-time collision risk, and introduce energy-aware altitude regulation that shapes the optimization preference for positive altitude gains."

"An energy-aware NMPC cost with an asymmetric climb penalty is introduced for 3D UAV obstacle avoidance. The additional penalty on positive altitude gains provides energy-aware altitude regulation by shaping the optimization preference among altitude motion, safety, and energy-related cost."

"To implement energy-aware altitude regulation in the local optimization, the energy-aware cost $J_{energy}$ is defined as"

"Here, $\phi_{\epsilon}(s)=\frac{1}{2}\left(s+\sqrt{s^2+\epsilon}\right)$ is a differentiable smooth approximation of the positive part of $s$, where $\epsilon>0$. The non-negative coefficients $c_1$, $c_2$, and $c_3$ weight the squared displacement, control-effort, and additional positive-altitude-gain terms, respectively. The third term assigns extra cost to positive altitude increments and therefore shapes the altitude-related optimization preference; it provides energy-aware altitude regulation rather than a guarantee of statistically significant mission-energy reduction. Thus, $J_{energy}$ remains a lightweight surrogate for online NMPC rather than a full rotor-level power model."

"This quantity is a model-based mission-energy proxy for comparative analysis, rather than a hardware battery measurement, a power measurement, or a flight-energy validation result."

"Table III reports the quantitative results over 30 randomized trials. The proposed ET-NMPC achieves a 100.0\% success rate, which is the highest among all methods. Compared with C-NMPC, it reduces the trigger rate from 100.0\% to 77.7\% and reduces the average CPU time from 2588.46 ms to 2163.62 ms. It also decreases the control effort from 119.83 to 107.92 while maintaining a comparable estimated energy proxy of 0.8082 Wh. Although NP-ET-NMPC obtains lower CPU time and a lower model-based energy proxy, its success rate drops to 83.3\%, and its clearance decreases to 1.44 m. Therefore, the proposed ET-NMPC achieves an overall trade-off among reliability, safety, computational efficiency, and energy-aware performance; the isolated influence of the asymmetric climb penalty is evaluated by the paired ablation in Table IV."

"The asymmetric climb and energy-related weights regulate the tradeoff between altitude motion and energy-related cost. These weights are fixed in all reported experiments; their isolated influence is evaluated by the paired ablation, which does not claim statistically significant reduction in cumulative positive climb or estimated mission energy."

"The estimated energy proxy of the proposed ET-NMPC remains comparable to the other NMPC-based methods. Therefore, these obstacle-count experiments characterize the overall robustness, safety margin, computational efficiency, and model-based energy-related behavior of the proposed framework; they do not isolate a statistically significant energy reduction attributable to the asymmetric climb penalty."

"The asymmetric climb penalty introduces an energy-aware altitude regulation mechanism by assigning additional cost to positive altitude gains. The paired ablation experiment evaluates its isolated influence without claiming statistically significant reduction in cumulative positive climb or estimated mission energy."

\begin{table}[h!]
\centering
\captionsetup{labelformat=empty}
\caption{\textcolor{blue}{Table IV: Paired Ablation Results for the Asymmetric Climb Penalty over 30 Matched Scenarios}}
\scriptsize
\setlength{\tabcolsep}{1.3pt}
\renewcommand{\arraystretch}{1.12}
\resizebox{\textwidth}{!}{%
\begin{tabular}{l c c c c c c c c c}
\toprule
\textcolor{blue}{\textbf{Algorithm}} & \textcolor{blue}{\textbf{Success}} & \textcolor{blue}{\textbf{Collision}} & \textcolor{blue}{\textbf{Total positive climb (m)}} & \textcolor{blue}{\textbf{Estimated energy proxy (Wh)}} & \textcolor{blue}{\textbf{Path length (m)}} & \textcolor{blue}{\textbf{Flight time (s)}} & \textcolor{blue}{\textbf{Minimum clearance (m)}} & \textcolor{blue}{\textbf{Trigger rate (\%)}} & \textcolor{blue}{\textbf{Control effort}} \\
\midrule
\textcolor{blue}{Asymmetric ET-NMPC} & \textcolor{blue}{30/30} & \textcolor{blue}{0/30} & \textcolor{blue}{11.52 $\pm$ 1.20} & \textcolor{blue}{0.22554 $\pm$ 0.01187} & \textcolor{blue}{60.55 $\pm$ 1.43} & \textcolor{blue}{13.55 $\pm$ 0.58} & \textcolor{blue}{1.93 $\pm$ 0.66} & \textcolor{blue}{77.92 $\pm$ 1.45} & \textcolor{blue}{107.51 $\pm$ 5.43} \\
\textcolor{blue}{Symmetric-energy ET-NMPC} & \textcolor{blue}{29/30} & \textcolor{blue}{1/30} & \textcolor{blue}{11.30 $\pm$ 1.23} & \textcolor{blue}{0.22186 $\pm$ 0.01807} & \textcolor{blue}{59.67 $\pm$ 4.60} & \textcolor{blue}{13.34 $\pm$ 1.09} & \textcolor{blue}{1.93 $\pm$ 0.66} & \textcolor{blue}{77.45 $\pm$ 2.69} & \textcolor{blue}{105.83 $\pm$ 9.60} \\
\bottomrule
\end{tabular}%
}
\begin{flushleft}
\textcolor{blue}{\textit{Note: Continuous metrics are reported as mean $\pm$ standard deviation over 30 trials. The paired differences are defined as Asymmetric $-$ Symmetric. Estimated energy is a model-based mission-energy proxy rather than a hardware battery measurement.}}
\end{flushleft}
\end{table}

The corresponding modifications have been highlighted in blue in the revised manuscript.

\medskip

\textcolor[rgb]{0.00,0.00,1.00}{[\textbf {Reviewer}]} \textit{\textbf{5. The rationale and representativeness of the baseline methods should be better justified.}}

\textit{\textbf{The manuscript compares the proposed method with continuous NMPC, two internal variants, 3D-APF, and RRT. These comparisons are useful, particularly for evaluating the effects of event triggering, prediction, and the energy formulation. However, continuous NMPC and the two event-triggered variants are primarily ablation baselines, whereas conventional APF and RRT may not represent strong recent competitors for prediction-based dynamic-obstacle avoidance.}}

\textit{\textbf{Several aspects of the comparison therefore require clarification. In particular, the manuscript should explain how the parameters of the baseline methods were tuned and whether all methods were evaluated under identical initial conditions, obstacle realizations, and random seeds. It should also clarify whether all methods received the same obstacle observations and, where applicable, whether the same motion-prediction information was used.}}

\textit{\textbf{The implementation of RRT* deserves particular attention. The manuscript should explain how RRT* was adapted to moving obstacles, how frequently replanning was performed, and how dynamic obstacle motion was incorporated into collision checking and closed-loop execution.}}

\textit{\textbf{More generally, the manuscript should justify why the selected baseline set is sufficient to establish competitiveness for online UAV navigation in dynamic environments. A comparison with, or at least a substantive discussion of, a recent prediction-based dynamic-obstacle avoidance method would make the experimental evaluation more convincing.}}

\textcolor[rgb]{0.00,0.00,1.00}{[\textbf{Authors}]}

Thank you very much for this valuable comment. We agree that the original manuscript did not sufficiently explain the roles and representativeness of the selected baselines, or make the comparison-fairness conditions sufficiently explicit.

In response, we have revised Section IV-B, Pages 7--8, to clarify the baseline-selection rationale. C-NMPC, NP-ET-NMPC, and SE-ET-NMPC are identified as internal ablation baselines for examining event-triggered activation, obstacle prediction, and the energy-aware formulation, respectively. In addition, 3D-APF and RRT* are identified as representative reactive-avoidance and sampling-based-planning methods.

We have also clarified that all methods are evaluated under identical initial conditions, static maps, dynamic-obstacle realizations, and random seeds. The NMPC-based variants share the same system model, prediction horizon, constraints, and solver settings. The revised text further specifies the common available obstacle observations, the distinct role of future obstacle-state estimation in prediction-based NMPC variants, and the closed-loop RRT* implementation with periodic replanning, updated obstacle states, and collision checking during execution.

Finally, we added a brief discussion explaining that recent prediction-based dynamic-obstacle avoidance methods mainly emphasize motion prediction and trajectory generation, whereas the present work focuses on resource-aware online control and adaptive NMPC activation. These directions are therefore complementary rather than directly equivalent baselines.

The corresponding revisions have been made in Section IV-B, Pages 7--8. The relevant revised text is reproduced below:

"C-NMPC, NP-ET-NMPC, and SE-ET-NMPC serve as internal ablation baselines for examining event-triggered activation, obstacle prediction, and the energy-aware formulation, respectively."

"3D-APF and RRT* represent reactive-avoidance and sampling-based-planning paradigms. Thus, the baseline set spans predictive control, event-triggered operation, energy-aware optimization, reactive planning, and sampling-based planning."

"All algorithms are compared under identical initial conditions, static maps, dynamic-obstacle realizations, and random seeds. The NMPC-based variants share the same system model, prediction horizon, constraints, and solver settings."

"All methods use the same available obstacle observations, while only the prediction-based NMPC variants exploit future obstacle-state estimation within their original formulations."

"RRT* is executed in closed loop through periodic replanning with updated obstacle states and collision checking during execution."

"RRT* shows limited performance under dynamic obstacle motion because it does not explicitly optimize future obstacle interactions during closed-loop execution."

"Recent prediction-based avoidance methods mainly address motion prediction and trajectory generation; they are complementary to, rather than directly equivalent to, the present focus on resource-aware online control and adaptive NMPC activation."

The corresponding modifications have been highlighted in blue in the revised manuscript.

\medskip

\textit{\textbf{Minor Comments}}

\textcolor[rgb]{0.00,0.00,1.00}{[\textbf {Reviewer}]} \textit{\textbf{1. Section III could be streamlined for improved readability.}}

\textit{\textbf{Section III is rather dense and combines obstacle prediction, risk assessment, event-triggered switching, NMPC formulation, theoretical discussion, solver implementation, and the complete algorithm in a single section. In addition, the explanation of solver activation and computational savings is repeated in several subsections.}}

\textit{\textbf{I suggest keeping the core methodological formulation in the main text while shortening repeated explanations and separating implementation details from the main algorithmic development. Some solver-specific or extended analytical details could also be moved to an appendix if appropriate.}}

\textcolor[rgb]{0.00,0.00,1.00}{[\textbf{Authors}]}

Thank you very much for this valuable comment. We agree that the original Section III contained a high density of methodological material because it combined obstacle prediction, risk assessment, event-triggered switching, NMPC formulation, theoretical analysis, solver implementation, and the complete algorithm. We also agree that some explanations of solver activation and computational savings were repeated across the section.

In response, we have streamlined Section III while retaining the core methodological formulas, theoretical analysis, and implementation transparency in the main text. The overall architecture is now focused on the observation--prediction--risk-evaluation--switching--control pipeline rather than performance conclusions. The risk-assessment subsection retains the TCPA/DCPA metrics, composite risk, and triggering conditions, while repeated descriptions of computational savings have been removed or condensed. In the theoretical discussion, experimental interpretations of CPU-time and Monte Carlo results have been removed so that the discussion remains focused on Zeno-free switching, computational complexity, recursive feasibility, and practical stability.

The solver description in Section III-F has been shortened from a detailed implementation tutorial to a concise implementation summary. It retains the use of multiple shooting, CasADi/IPOPT, a shifted-reference warm-start during triggered replanning, and lightweight tracking control in low-risk mode. Solver configurations and the expanded computational profiling disclosure remain available in Section IV-E, whereas Section III-F now only describes the method implementation needed to interpret the framework and Algorithm 1.

The corresponding revisions are shown below:

"This dual-mode design enables risk-dependent switching between lightweight tracking and predictive replanning according to the estimated collision risk."

"The NMPC optimization problem is implemented as a sparse nonlinear programming problem using multiple shooting. The resulting NLP is solved through the CasADi/IPOPT framework. During triggered replanning, a shifted-reference warm-start strategy is adopted, while low-risk periods bypass the optimization module and apply a lightweight tracking controller. Detailed solver configurations and computational profiling results are provided in Section IV-E."

"This section provides theoretical analysis for the proposed event-triggered dual-mode NMPC framework. The analysis first proves that the triggering mechanism excludes Zeno-like mode chattering. It then characterizes the computational complexity of prediction and event-triggered solver activation. Finally, recursive feasibility and closed-loop stability are discussed under bounded prediction errors and feasible replanning conditions."

"The complete procedure is summarized in Algorithm 1."

The corresponding modifications have been highlighted in blue in the revised manuscript.

\medskip

\textcolor[rgb]{0.00,0.00,1.00}{[\textbf {Reviewer}]} \textit{\textbf{2. Table III should report statistical variation and provide more information about the randomized trials.}}

\textit{\textbf{The manuscript states that Table III is based on 30 randomized trials, but only point estimates are reported. Standard deviations, confidence intervals, or other measures of statistical variation should be provided for the main performance metrics.}}

\textit{\textbf{It should also be clarified whether the algorithms were evaluated using paired random seeds and exactly the same obstacle trajectories in each trial. This is particularly important because several differences in clearance, climb, and estimated energy consumption are relatively small.}}

\textit{\textbf{For such comparisons, confidence intervals for paired differences, or another appropriate paired statistical analysis, would help determine whether the observed differences are practically meaningful.}}

\textcolor[rgb]{0.00,0.00,1.00}{[\textbf{Authors}]}

Thank you for this important comment. We agree that the original presentation of Table III did not provide sufficient statistical variation or explicitly describe whether the comparisons were paired. In particular, point estimates alone were not sufficient to determine whether the small differences in climb and estimated energy were statistically meaningful.

In response, we have added a 30-trial paired Monte Carlo ablation in Section IV-C and Table IV. The study isolates the asymmetric positive-climb penalty by comparing Asymmetric ET-NMPC and Symmetric-energy ET-NMPC under matched scenarios. It reports the continuous metrics as mean $\pm$ standard deviation and evaluates paired differences using 95\% confidence intervals. Table III remains the overall comparison of heterogeneous baseline methods, whereas the new paired analysis provides the appropriate statistical evidence for the isolated energy-aware formulation comparison.

To further improve the statistical transparency of the main Table III comparison, we also conducted a paired-difference analysis between ET-NMPC and C-NMPC using the 30 matched Monte Carlo trials. Both methods used the same scenario seed, initial condition, and dynamic-obstacle realization in each matched trial. Because C-NMPC failed in trial 2, continuous metrics were evaluated over the 29 trials with valid values for both methods. The mean CPU-time difference (C-NMPC $-$ ET-NMPC) was 425.11 ms, with a 95\% confidence interval of [372.62, 477.60] ms. The paired minimum-clearance difference (ET-NMPC $-$ C-NMPC) was 0.009 m (95\% CI: [$-0.076$, $0.094$] m), and the paired control-effort difference (C-NMPC $-$ ET-NMPC) was 11.83 (95\% CI: [10.82, 12.84]). These paired estimates provide statistical support for the reported computational advantage, while the clearance interval includes zero and therefore does not indicate a systematic degradation of the safety margin in the matched valid trials.

The corresponding revisions are shown below:

\textcolor{blue}{%
\begin{table}[h!]
\color{blue}
\centering
\captionsetup{labelformat=empty}
\caption{Table III: Quantitative Performance Comparison of Different Algorithms in a Complex Dynamic Urban Scenario}
\scriptsize
\renewcommand{\arraystretch}{1.12}
\setlength{\tabcolsep}{1.5pt}
\resizebox{\textwidth}{!}{%
\begin{tabular}{l c c c c c c c}
\toprule[1.1pt]
\textbf{Algorithm} & \textbf{Success (\%)} & \textbf{Trigger (\%)} & \textbf{Avg CPU (ms)} & \textbf{Clearance (m)} & \textbf{Effort} & \textbf{Climb (m)} & \textbf{Est. energy proxy (Wh)} \\
\midrule
\textbf{Proposed ET-NMPC} & 100.0 (30/30) & 77.7 $\pm$ 1.4 & 2163.62 $\pm$ 70.39 & 1.94 $\pm$ 0.67 & 107.92 $\pm$ 5.10 & 11.59 $\pm$ 1.17 & 0.8082 $\pm$ 0.0336 \\
C-NMPC & 96.7 (29/30) & 100.0 $\pm$ 0.0 & 2588.46 $\pm$ 136.94 & 1.95 $\pm$ 0.63 & 119.83 $\pm$ 5.36 & 11.67 $\pm$ 1.32 & 0.8199 $\pm$ 0.0389 \\
NP-ET-NMPC & 83.3 (25/30) & 75.0 $\pm$ 0.4 & 1862.95 $\pm$ 11.26 & 1.44 $\pm$ 0.86 & 103.56 $\pm$ 1.47 & 10.97 $\pm$ 0.38 & 0.7944 $\pm$ 0.0078 \\
SE-ET-NMPC & 96.7 (29/30) & 77.7 $\pm$ 1.4 & 1918.71 $\pm$ 35.32 & 1.96 $\pm$ 0.67 & 107.65 $\pm$ 4.85 & 11.46 $\pm$ 1.17 & 0.8047 $\pm$ 0.0310 \\
3D-APF & 13.3 (4/30) & 0.0 $\pm$ 0.0 & 0.13 $\pm$ 0.12 & 0.54 $\pm$ 0.43 & 132.03 $\pm$ 15.67 & 11.17 $\pm$ 1.34 & 0.9548 $\pm$ 0.1449 \\
RRT* & 0.0 (0/30) & -- & -- & -- & -- & -- & -- \\
\bottomrule[1.1pt]
\end{tabular}
}
\begin{flushleft}
\footnotesize
\textit{Note: Success is reported as rate (successful trials/30). Collision counts for Proposed ET-NMPC, C-NMPC, NP-ET-NMPC, SE-ET-NMPC, 3D-APF, and RRT* are 0/30, 1/30, 5/30, 1/30, 26/30, and 30/30, respectively. Continuous metrics are reported as mean $\pm$ standard deviation over successful trials. For failed trials, only success and collision statistics are counted.}
\end{flushleft}
\end{table}
}

"To isolate the influence of the asymmetric climb penalty, we conducted 30 paired Monte Carlo scenarios comparing Asymmetric ET-NMPC with Symmetric-energy ET-NMPC. For each paired scenario, both controllers used the same random seed, static map, UAV initial state, goal position, dynamic-obstacle realization, and measurement-noise sequence. The only changed setting was the additional positive-climb penalty: the asymmetric controller used $c_3=3.0$, whereas the symmetric-energy controller disabled this term by setting $c_3=0$."

"Note: Continuous metrics are reported as mean $\pm$ standard deviation over 30 trials. Paired differences are defined as Asymmetric $-$ Symmetric. Estimated energy is a model-based mission-energy proxy rather than a hardware battery measurement."

"The asymmetric controller achieved 30/30 successful trials with no collision, whereas the symmetric-energy controller achieved 29/30 successful trials with one collision. The paired difference in total positive climb was 0.217 m (95\% CI: [$-0.069$, $0.503$]), and the paired difference in estimated energy proxy was 0.00368 Wh (95\% CI: [$-0.00169$, $0.00904$]), where each difference is defined as Asymmetric $-$ Symmetric. Both confidence intervals include zero. Therefore, this experiment does not provide statistical evidence that the asymmetric penalty reduces cumulative positive climb or estimated mission energy under the tested parameterization. Instead, the additional term shapes altitude-related optimization preference and provides energy-aware altitude regulation within the safety--energy tradeoff."

The corresponding modifications have been highlighted in blue in the revised manuscript.

%%%%%%%%%%%%%%%%%%%%%%%%%%%%%%%%%%%%%%%%%%%%%%%%%%%%%%%%%%%%%%
\begin{center}
\textbf{\textit{\large Response to Reviewer \# 2}}
\end{center}

\textcolor[rgb]{0.00,0.00,1.00}{[\textbf {Reviewer}]} \textit{\textbf{Comments to the Author}}

\textit{\textbf{I have carefully reviewed the revised manuscript, along with the authors' responses to all reviewer comments. In my assessment, the authors have adequately addressed the concerns raised and have made substantial and appropriate revisions to the text. Overall, the manuscript now meets the publication standards of this journal. As a minor suggestion, I note that the reference list includes relatively few papers published within the last two years; updating some of the references with more recent literature would further enhance the manuscript's timeliness and relevance.}}

\textcolor[rgb]{0.00,0.00,1.00}{[\textbf{Authors}]}

%%%%%%%%%%%%%%%%%%%%%%%%%%%%%%%%%%%%%%%%%%%%%%%%%%%%%%%%%%%%%%
\begin{center}
\textbf{\textit{\large Response to Reviewer \# 3}}
\end{center}

\textcolor[rgb]{0.00,0.00,1.00}{[\textbf {Reviewer}]} \textit{\textbf{Comments to the Author}}

\textit{\textbf{I have no more comments for this manuscript.}}

\textcolor[rgb]{0.00,0.00,1.00}{[\textbf{Authors}]}

\end{document}
